\chapter{General Conclusion}
\chaptertoc
\clearpage

\section{Summary}
This end-of-studies project produced a queue wait-time estimation platform that connects computer vision, queue analytics, web supervision, and alert automation. The project started from a practical problem: managers need more than a visual impression of a queue. They need indicators that show how the queue is evolving and when action is required.

In practical terms, the system now includes:
\begin{itemize}
    \item a perception pipeline based on YOLO detection and ByteTrack identity continuity,
    \item queue-zone filtering and event-driven estimation of $\lambda$, $\mu$, and expected wait time,
    \item adaptive wait-time outputs that respond to queue conditions,
    \item a secure authentication layer for manager sign-up, sign-in, session refresh, and administrator account governance,
    \item a typed FastAPI backend-for-frontend with feed lifecycle control and live websocket events,
    \item a React dashboard for monitoring, setup, and runtime interaction,
    \item n8n and Telegram integration for operational alert routing.
\end{itemize}

Taken together, these elements form a monitoring-to-action loop: the system observes a queue, estimates its state, displays it to the manager, and sends alerts when the situation requires attention.

Across the four sprint chapters, the work progressed from perception and queue analytics to live services, then to the web dashboard, and finally to automation and validation.

\section{Objective Review}
The initial objectives were achieved at the project scope level:
\begin{enumerate}
    \item \textbf{Real-time detection and tracking:} achieved through YOLO-based detection and ByteTrack identity continuity.
    \item \textbf{Queue-zone counting:} achieved with polygon-zone management and in-zone identity filtering.
    \item \textbf{Queue metric estimation:} achieved for arrival rate, service rate, and wait time.
    \item \textbf{Adaptive wait-time estimation:} achieved through automatic formula switching based on queue load.
    \item \textbf{Web monitoring:} achieved through a multi-feed dashboard with live updates, setup workflows, and fallback behavior.
    \item \textbf{Automation integration:} achieved through webhook transport, n8n routing, and Telegram notifications.
\end{enumerate}

The main remaining gap is not architectural completeness, but the depth of long-duration evaluation under broader operational scenarios.

\section{Future Improvements}
Several extensions would strengthen the platform before wider deployment:
\begin{itemize}
    \item \textbf{Mobile application}: developing a companion mobile application allowing operators and managers to monitor queue activity, consult statistics, and receive alerts remotely in real time.
    \item \textbf{In-platform messaging}: introducing a notification system enabling managers to communicate directly with operators, such as assigning available staff to queues based on live crowd conditions.
    \item \textbf{Tracking accuracy improvement}: refining the person re-identification pipeline to reduce identity switches and tracking losses, particularly in dense or partially occluded queue scenarios.
\end{itemize}

